%=========================================================================
%  ProcGuard — Informe técnico final
%  Curso de Sistemas Operativos y Laboratorio
%  Universidad de Antioquia — Facultad de Ingeniería — 2026-1
%
%  Compilar en Overleaf con la clase IEEEtran (pdfLaTeX).
%=========================================================================
\documentclass[conference]{IEEEtran}
\IEEEoverridecommandlockouts

\usepackage[utf8]{inputenc}
\usepackage[T1]{fontenc}
\usepackage[spanish]{babel}
\usepackage{graphicx}
\usepackage{amsmath}
\usepackage{booktabs}
\usepackage{url}
\usepackage{listings}
\usepackage{xcolor}

\lstset{
  basicstyle=\ttfamily\footnotesize,
  breaklines=true,
  frame=single,
  columns=fullflexible,
  keepspaces=true
}

\begin{document}

\title{ProcGuard: Sistema de Gobernanza Activa de Procesos para Linux\\
\large Monitoreo multidimensional y respuesta correctiva escalonada sobre procfs y cgroups~v2}

\author{
\IEEEauthorblockN{Brahian Ocampo García}
\IEEEauthorblockA{Facultad de Ingeniería\\
Universidad de Antioquia\\
Curso de Sistemas Operativos y Laboratorio, 2026-1\\
\texttt{brahian.ocampo@udea.edu.co}}
}

\maketitle

%-------------------------------------------------------------------------
\begin{abstract}
Las herramientas tradicionales de monitoreo de procesos en Linux (\texttt{top},
\texttt{htop}) son pasivas: muestran métricas pero desconectan la observación de
la acción correctiva, que queda como una secuencia manual de comandos sujeta a
latencia, error humano y ausencia de trazabilidad. Este trabajo presenta
\emph{ProcGuard}, un sistema escrito en C~(C11) que implementa el ciclo
observación~$\rightarrow$~detección~$\rightarrow$~decisión~$\rightarrow$~acción
sobre los procesos en ejecución, gobernado por políticas declarativas en formato
INI. A partir de \texttt{procfs} se calculan métricas derivadas (uso de CPU,
memoria residente, tasas de E/S) y se evalúan contra umbrales con tres
mecanismos contra falsos positivos: persistencia temporal, histéresis y
\emph{cooldown}. Las anomalías persistentes desencadenan un escalamiento
progresivo de acciones correctivas
(\texttt{warn}~$\rightarrow$~\texttt{renice}~$\rightarrow$~\texttt{cage}~$\rightarrow$~\texttt{stop}~$\rightarrow$~\texttt{kill})
con protecciones contra daño autoinfligido (lista blanca, techo de acciones,
modo \emph{dry-run} y revalidación de identidad \texttt{(pid, starttime)}). El
confinamiento se materializa mediante \texttt{cgroups~v2}. Se entrega el núcleo
de gobernanza (módulos M1--M4) verificado con 111 pruebas unitarias bajo
AddressSanitizer y UndefinedBehaviorSanitizer. Se reporta la metodología, la
arquitectura y un diseño experimental formal para evaluación estadística de
precisión de detección y efectividad de gobernanza, junto con los resultados de
verificación obtenidos.
\end{abstract}

\begin{IEEEkeywords}
sistemas operativos, procfs, cgroups, gobernanza de procesos, monitoreo,
control de recursos, Linux.
\end{IEEEkeywords}

%-------------------------------------------------------------------------
\section{Introducción}

Los sistemas Linux modernos operan bajo cargas mixtas de CPU, memoria, E/S y red,
con servicios, contenedores y máquinas virtuales coexistiendo sobre el mismo
hardware. En este contexto, la observación pasiva resulta insuficiente: un
administrador que detecta un proceso problemático mediante \texttt{top} debe
identificarlo, evaluar su impacto, decidir una acción y ejecutarla mediante
comandos separados (\texttt{kill}, \texttt{renice}, \texttt{cgcreate}). Este
flujo manual introduce latencia, riesgo de error y falta de trazabilidad.

Linux ofrece a nivel de kernel mecanismos de medición y control
---\texttt{procfs}, \texttt{cgroups~v2}, \texttt{setpriority()},
\texttt{sched\_setaffinity()}--- que permiten observar y actuar
programáticamente sobre los procesos. Sin embargo, ninguna herramienta estándar
de terminal integra en un solo sistema la observación multidimensional, la
detección de anomalías con control de falsos positivos y la respuesta
automatizada con protecciones de seguridad.

\textbf{Contribución.} Este trabajo diseña e implementa el núcleo de
\emph{ProcGuard}, un sistema de gobernanza activa que cierra el lazo entre
observación y acción mediante políticas configurables. El objetivo general es
integrar el ciclo observación, detección, decisión y acción sobre procesos,
gobernado por políticas definidas por el administrador. Como objetivos
específicos del alcance entregado: (i)~implementar un recolector basado en
\texttt{procfs} con manejo resiliente de errores; (ii)~calcular métricas
derivadas a partir de muestras consecutivas; (iii)~desarrollar un motor de
detección con persistencia, histéresis y \emph{cooldown}; y (iv)~implementar un
motor de gobernanza con escalamiento progresivo y protecciones contra daño
autoinfligido.

%-------------------------------------------------------------------------
\section{Marco teórico}

\subsection{El pseudo-sistema de archivos \texttt{procfs}}
\texttt{procfs} es la interfaz primaria del kernel para exponer información del
sistema y de los procesos mediante archivos virtuales legibles. ProcGuard usa
\texttt{/proc/[pid]/stat} (estado, tiempos de CPU, prioridad, \texttt{starttime}),
\texttt{/proc/[pid]/status} (memoria \texttt{VmRSS}/\texttt{VmSize}) y
\texttt{/proc/[pid]/io} (bytes leídos/escritos). La tupla
\texttt{(pid, starttime)} se emplea como identificador único para evitar errores
por reciclaje de PID.

\subsection{Cálculo de métricas derivadas}
El uso de CPU de un proceso no es un valor instantáneo, sino una tasa derivada
entre dos muestras consecutivas. Sea $\Delta t_{proc} = (utime+stime)$ la
variación de \emph{ticks} de CPU del proceso y $\Delta t_{wall}$ el intervalo de
tiempo de pared medido con \texttt{clock\_gettime(CLOCK\_MONOTONIC)}. Con $\mathit{HZ}$
ticks por segundo y $n_{cpu}$ núcleos:
\begin{equation}
\%\mathrm{CPU} = \frac{\Delta t_{proc} / \mathit{HZ}}{\Delta t_{wall}\cdot n_{cpu}}\times 100.
\end{equation}
El uso de \texttt{CLOCK\_MONOTONIC} garantiza robustez frente a ajustes del reloj
del sistema. La memoria se mide con \texttt{VmRSS} (Resident Set Size), de bajo
costo de lectura y suficiente para la detección de anomalías de consumo.

\subsection{Mecanismos de control de recursos del kernel}
Las acciones correctivas se apoyan en cuatro mecanismos: señales POSIX
(\texttt{SIGSTOP}, \texttt{SIGCONT}, \texttt{SIGTERM}, \texttt{SIGKILL}); la
prioridad \emph{nice} vía \texttt{setpriority()}; la afinidad de CPU vía
\texttt{sched\_setaffinity()}; y \texttt{cgroups~v2}, que permite confinar un
proceso bajo \texttt{/sys/fs/cgroup/} aplicando límites de CPU (\texttt{cpu.max})
y memoria (\texttt{memory.max}) sin terminarlo.

%-------------------------------------------------------------------------
\section{Metodología}

El desarrollo siguió un enfoque incremental dirigido por pruebas (TDD): cada
módulo se construyó mediante ciclos RED--GREEN, escribiendo primero las pruebas
que especifican el comportamiento esperado y luego la implementación mínima que
las satisface. El trabajo se organizó en \emph{slices} verticales con
verificación continua bajo sanitizers de memoria y comportamiento indefinido.

La estrategia de pruebas se estructura en tres niveles: (i)~pruebas unitarias por
módulo con el framework Unity; (ii)~pruebas de integración con un \texttt{procfs}
sintético, posible porque la ruta base de \texttt{/proc} es un parámetro
configurable (\texttt{-{}-proc}), lo que permite construir un árbol controlado en
\texttt{/tmp} y validar el \emph{pipeline} de forma determinista; y
(iii)~pruebas extremo a extremo del binario sobre el sistema real. Todo el código
del proyecto se compila con
\texttt{-Wall -Wextra -Werror -Wshadow} y se ejecuta bajo
\texttt{-fsanitize=address,undefined} durante las pruebas.

%-------------------------------------------------------------------------
\section{Implementación}

ProcGuard se diseñó como una arquitectura de siete módulos
(Fig.~\ref{fig:arch}). Este informe documenta el núcleo de gobernanza
\textbf{implementado y verificado (M1--M4)} y delimita explícitamente el alcance
restante como trabajo futuro (Sec.~\ref{sec:alcance}).

\begin{figure}[t]
\centering
\includegraphics[width=\columnwidth]{arquitectura.jpg}
\caption{Arquitectura de ProcGuard. El \emph{pipeline} de datos M1--M3 alimenta
el motor de gobernanza M4. Los módulos M5--M7 quedan especificados como trabajo
futuro.}
\label{fig:arch}
\end{figure}

\subsection{M1 — Recolector de datos}
Responsable de la lectura directa de \texttt{procfs}. En cada ciclo escanea el
directorio base para descubrir procesos activos y almacena, por proceso, la tupla
\texttt{(pid, starttime)}. Aplica una política de resiliencia: los fallos
\texttt{ENOENT}, \texttt{ESRCH} y \texttt{EACCES} se tratan como condiciones
normales y se descartan silenciosamente; el recolector nunca aborta por un
proceso individual que falla.

\subsection{M2 — Almacén de muestras}
Buffer circular de las últimas $N$ muestras por proceso. Los procesos que
desaparecen se conservan durante una ventana de gracia de $G$ ciclos para
permitir que las evaluaciones pendientes se completen, tras lo cual sus
estructuras se liberan.

\subsection{M3 — Motor de métricas}
Transforma muestras crudas en métricas accionables: \texttt{cpu\_percent},
memoria residente y virtual, y tasas de E/S (\texttt{io\_read\_rate},
\texttt{io\_write\_rate}) calculadas como deltas entre muestras consecutivas.

\subsection{M4 — Motor de alertas y gobernanza}
Es el núcleo diferenciador. Evalúa las métricas contra políticas declarativas e
implementa tres mecanismos contra falsos positivos:
\begin{itemize}
\item \textbf{Persistencia temporal:} la alerta se activa solo si el umbral se
supera durante $N$ muestras consecutivas.
\item \textbf{Histéresis:} la desactivación requiere que la métrica descienda por
debajo de un umbral inferior durante $M$ muestras.
\item \textbf{Cooldown:} tras cada acción correctiva se espera un período de
enfriamiento antes de escalar.
\end{itemize}

El motor soporta \emph{escalamiento progresivo}: una secuencia ordenada de
acciones de severidad creciente
(\texttt{warn}~$\rightarrow$~\texttt{renice}~$\rightarrow$~\texttt{cage}~$\rightarrow$~\texttt{stop}~$\rightarrow$~\texttt{kill}).
El confinamiento (\texttt{cage}) crea un \texttt{cgroup~v2} real con límites de
CPU. Las políticas se definen en formato INI extendido, parseado con la librería
embebida \texttt{inih}:

\begin{lstlisting}
[policy:cpu_hog]
type = performance
metric = cpu_percent
threshold = 80.0
threshold_low = 60.0
persistence = 3
hysteresis_m = 2
cooldown_s = 10
actions = warn, renice:10, cage, stop, kill
\end{lstlisting}

\begin{figure}[t]
\centering
\includegraphics[width=\columnwidth]{ciclo_gobernanza.jpg}
\caption{Ciclo de gobernanza. Una métrica que supera el umbral debe persistir
$N$ muestras para activar la alerta; tras pasar las protecciones del sistema, se
ejecuta el escalamiento progresivo, regulado por histéresis y \emph{cooldown}.}
\label{fig:ciclo}
\end{figure}

El flujo completo de decisión se resume en la Fig.~\ref{fig:ciclo}.

\subsection{Protecciones contra daño autoinfligido}
Las protecciones no limitan la detección, solo la acción automática:
(i)~\emph{lista blanca} de procesos protegidos (PID~1, el propio ProcGuard,
procesos del kernel y nombres protegidos que residan en rutas estándar);
(ii)~\emph{techo de acciones} (máximo de \emph{kills} por minuto y de procesos
confinados simultáneamente); (iii)~\emph{modo dry-run}, activo por defecto, que
ejecuta todo el ciclo pero registra las acciones en lugar de ejecutarlas; y
(iv)~\emph{revalidación de identidad}: antes de actuar se verifica que
\texttt{(pid, starttime)} no cambió, previniendo actuar sobre un PID reciclado
(condición de carrera TOCTOU).

%-------------------------------------------------------------------------
\section{Protocolo de experimentación}

Se formuló un diseño experimental para dar certeza estadística a la evaluación de
ProcGuard, estructurado en dos dimensiones, con un mínimo de 12 repeticiones por
tratamiento y orden aleatorizado.

\textbf{Dimensión 1 — Funcionalidad.} Evalúa la precisión de detección y la
efectividad de gobernanza. Factores: tipo de carga (CPU-bound, memory-bound),
perfil de anomalía (pico instantáneo, sostenido) y política (solo alertas,
escalamiento completo). Variables de respuesta: latencia de detección (ms), tasa
de verdaderos positivos, tasa de falsos positivos y efectividad de gobernanza.

\textbf{Protocolo común por tratamiento.} Preparación a estado limpio; línea base
de 60~s sin el monitor; inicio del monitor y estabilización de 30~s; inyección de
la carga sintética con \texttt{stress-ng}; observación de 120~s registrando todas
las variables; recolección de \emph{logs}.

\textbf{Métodos estadísticos planeados.} Estadística descriptiva con intervalos
de confianza al 95\%; verificación de supuestos con Shapiro--Wilk (normalidad) y
Levene (homocedasticidad); ANOVA de una vía para comparar latencias entre
tratamientos; y Tukey HSD como prueba \emph{post-hoc} para identificar los pares
de tratamientos con diferencias significativas. El análisis se implementa en
Python (\texttt{scipy}).

%-------------------------------------------------------------------------
\section{Resultados}

Por restricciones de tiempo, la ejecución del experimento estadístico de la
sección anterior quedó pendiente; en consecuencia, los resultados aquí reportados
corresponden a la \emph{verificación y validación funcional} del núcleo
implementado (M1--M4).

\textbf{Verificación por pruebas unitarias.} El sistema se validó con 111 pruebas
unitarias distribuidas por módulo (Tabla~\ref{tab:tests}), todas en estado
\textsc{PASS}, ejecutadas bajo AddressSanitizer y UndefinedBehaviorSanitizer sin
reportar fugas de memoria ni comportamiento indefinido.

\begin{table}[t]
\caption{Cobertura de pruebas unitarias por módulo.}
\label{tab:tests}
\centering
\begin{tabular}{lc}
\toprule
\textbf{Suite de pruebas} & \textbf{N.º de pruebas} \\
\midrule
Motor de alertas (engine) & 58 \\
Recolector (collector) & 12 \\
Métricas (metrics) & 10 \\
Almacén (store) & 10 \\
Estado de alertas (alert\_state) & 8 \\
Parser de políticas (alert\_parser) & 7 \\
Colas IPC (queues) & 6 \\
\midrule
\textbf{Total} & \textbf{111} \\
\bottomrule
\end{tabular}
\end{table}

\textbf{Validación funcional.} El binario ejecuta el ciclo completo de gobernanza
sobre \texttt{/proc} e imprime, por ciclo, el ranking de procesos por \%CPU. Las
pruebas verifican comportamientos críticos de forma determinista, entre ellos:
la cancelación de la acción cuando \texttt{starttime} cambia (TOCTOU); la
supresión de llamadas al sistema en modo \emph{dry-run} con avance del nivel de
escalamiento; la aplicación idempotente del \emph{cage} con conteo correcto bajo
el techo de procesos confinados; el bloqueo de \texttt{kill} al alcanzar el techo
de acciones; y la liberación del confinamiento durante la recolección de basura
cuando el proceso desaparece. Estos resultados evidencian que las protecciones
contra daño autoinfligido operan según lo especificado antes de habilitar
acciones reales.

%-------------------------------------------------------------------------
\section{Alcance y trabajo futuro}
\label{sec:alcance}

El núcleo de gobernanza (M1--M4) está implementado, probado y operativo,
incluyendo el confinamiento real con \texttt{cgroups~v2}. Quedan especificados y
diseñados, como trabajo futuro, los siguientes componentes: el motor de seguridad
(M5) con cuatro heurísticas de comportamiento sospechoso (escaneo de puertos,
procesos camuflados, enumeración del sistema vía \texttt{inotify} y huérfanos
anómalos); la interfaz de terminal interactiva (M6) con \texttt{ncurses}; el
subsistema de reportes y registro forense (M7) con \emph{logs} JSON-lines y
reportes HTML; el modo \emph{daemon} con señales POSIX; el modelo de tres hilos
con sincronización por \emph{mutex}; las métricas de red e hilos; y la ejecución
del diseño experimental con su análisis estadístico.

%-------------------------------------------------------------------------
\section{Conclusiones}

Se diseñó e implementó el núcleo de un sistema de gobernanza activa de procesos
para Linux que cierra el lazo entre observación y acción, superando el paradigma
pasivo de las herramientas convencionales. El motor de gobernanza combina
detección con control de falsos positivos (persistencia, histéresis,
\emph{cooldown}) y respuesta correctiva escalonada respaldada por mecanismos
reales del kernel (\texttt{setpriority}, señales POSIX y \texttt{cgroups~v2}). El
énfasis en protecciones contra daño autoinfligido ---en particular la
revalidación de identidad \texttt{(pid, starttime)} y el modo \emph{dry-run} por
defecto--- resultó central para un sistema que actúa automáticamente sobre el
sistema operativo. La verificación con 111 pruebas unitarias bajo sanitizers
respalda la corrección del núcleo entregado. El principal trabajo pendiente es la
ejecución del diseño experimental formulado, que permitirá cuantificar
estadísticamente la latencia de detección y la efectividad de gobernanza, además
de completar los módulos de seguridad, interfaz y reportes ya especificados.

%-------------------------------------------------------------------------
\begin{thebibliography}{9}

\bibitem{procfs}
\textit{proc(5) — process information pseudo-filesystem}, Linux \emph{man-pages}.
\url{https://man7.org/linux/man-pages/man5/proc.5.html}

\bibitem{cgroups}
\textit{cgroups(7) — Linux control groups}, Linux \emph{man-pages}.
\url{https://man7.org/linux/man-pages/man7/cgroups.7.html}

\bibitem{sched}
\textit{sched\_setaffinity(2)} y \textit{setpriority(2)}, Linux \emph{man-pages}.
\url{https://man7.org/linux/man-pages/man2/setpriority.2.html}

\bibitem{inih}
B. Hoyt, \textit{inih: simple .INI file parser in C}.
\url{https://github.com/benhoyt/inih}

\bibitem{unity}
ThrowTheSwitch, \textit{Unity: unit testing for C}.
\url{https://github.com/ThrowTheSwitch/Unity}

\bibitem{stressng}
C. I. King, \textit{stress-ng: a tool to load and stress a computer system}.
\url{https://github.com/ColinIanKing/stress-ng}

\end{thebibliography}

\end{document}
